% !PRINTED 151
\weberpage{151}
%<<CONT
\[ \mathrm{cn}^2v = 1-\mathrm{sn}^2v, \qquad \mathrm{dn}^2v = 1-\varkappa^2\mathrm{sn}^2v, \]
\[
\begin{aligned}
\frac{\sigma_{01}^2(u)}{\sigma^2(u)}-\frac{\sigma_{10}^2(u)}{\sigma^2(u)} &= \frac{4K^2}{\omega_1^2}, \\
\frac{\sigma_{01}^2(u)}{\sigma^2(u)}-\frac{\sigma_{00}^2(u)}{\sigma^2(u)} &= \frac{4K^2}{\omega_1^2}\varkappa^2.
\end{aligned}
\]

Three quantities $e_1$, $e_2$, $e_3$ independent of $u$ (and therefore depending only on
$\omega_1$, $\omega_2$) can accordingly be determined in such a way that
\weq{3}{\frac{\sigma_{10}^2(u)}{\sigma^2(u)}+e_1 = \frac{\sigma_{00}^2(u)}{\sigma^2(u)}+e_2 = \frac{\sigma_{01}^2(u)}{\sigma^2(u)}+e_3 = \wp(u),}
where $\wp(u)$ is a doubly periodic function newly defined by (3). The
quantities $e_1$, $e_2$, $e_3$ may be arbitrary, provided only that they satisfy the conditions
\weq{4}{e_1-e_3 = \frac{4K^2}{\omega_1^2}, \quad e_2-e_3 = \frac{4K^2}{\omega_1^2}\varkappa^2}
We are therefore free to add, for their complete determination, the further condition
\weq{5}{e_1+e_2+e_3 = 0}
The following expressions then result for $e_1$, $e_2$, $e_3$:
\weq{6}{\begin{aligned}
e_1 &= \frac{4K^2}{\omega_1^2}\frac{1+\varkappa'^2}{3}, \\
e_2 &= -\frac{4K^2}{\omega_1^2}\frac{\varkappa'^2-\varkappa^2}{3}, \\
e_3 &= -\frac{4K^2}{\omega_1^2}\frac{1+\varkappa^2}{3}.
\end{aligned}}

The function $\wp(u)$, which possesses the periods $\omega_1$, $\omega_2$, is hereby expressed by means of (1):
\weq{7}{\wp(u) = \frac{4K^2}{\omega_1^2}\left(\frac{1}{\mathrm{sn}^2v}-\frac{1+\varkappa^2}{3}\right).}

The function $\wp(u)$ has the desired property of invariability under linear
transformations, as one recognises from the expression
\weq{8}{\wp(u) = \frac{1}{3}\frac{\sigma_{00}^2(u)+\sigma_{01}^2(u)+\sigma_{10}^2(u)}{\sigma^2(u)}}

In consequence of (3) the $e_1$, $e_2$, $e_3$ are permuted under a linear transformation
in the same way as the $\sigma_{10}$, $\sigma_{00}$, $\sigma_{01}$; a symmetric function
of them is therefore, under a linear trans-
%CONT>>
